%% 由 build/emit_tex.py 生成（生成物，勿手改）；定制请放 user_style.tex / user_meta.yaml
\chapter{线性代数与矩阵}\label{ux7ebfux6027ux4ee3ux6570ux4e0eux77e9ux9635}

\begin{quote}
本附录讲科学计算里最常用的线性代数：向量与矩阵是什么、怎么''想''它们、数字上怎么算、Python
里怎么调。学完你应能用''线性变换''的眼光看明白第 1 章 NumPy
的线性代数、第 8 章 sklearn 的降维。
\end{quote}

\begin{center}\rule{0.5\linewidth}{0.5pt}\end{center}

\section{A.1
直觉故事：矩阵到底在干什么}\label{a.1-ux76f4ux89c9ux6545ux4e8bux77e9ux9635ux5230ux5e95ux5728ux5e72ux4ec0ux4e48}

很多同学背熟了矩阵乘法，却不知道它''在算什么''。我们从一个具体问题开始：

\begin{quote}
一家店只卖两种套餐：A 套餐（米饭 1 份 + 鸡肉 2 份），B 套餐（米饭 2 份 +
鸡肉 1 份）。今天卖出 3 份 A、2 份 B，问一共用了多少米饭和鸡肉？
\end{quote}

把配额写成向量：卖出的数量 \texttt{x\ =\ (3,\ 2)}（A 份数、B
份数）。每种套餐的用料写成列向量：A 套餐 \texttt{(1,\ 2)}，B 套餐
\texttt{(2,\ 1)}。那么总用料就是\textbf{线性组合}：

\begin{Shaded}
\begin{Highlighting}[]
\NormalTok{总米饭 = 3×1 + 2×2 = 7；总鸡肉 = 3×2 + 2×1 = 8}
\NormalTok{→ 写成矩阵： [1 2; 2 1] × [3; 2] = [7; 8]}
\end{Highlighting}
\end{Shaded}

这个例子说明了三件事：

\begin{enumerate}
\def\labelenumi{\arabic{enumi}.}
\tightlist
\item
  \textbf{向量是''配比''}：一行数据、一份订单、一个状态，都可以是一个向量；
\item
  \textbf{矩阵是''线性变换的约定''}：\texttt{A}
  的每一列是一种''原料的配方''，\texttt{Ax} 就是按 \texttt{x}
  的比例把各列混合起来；
\item
  \textbf{矩阵乘法 = 复合变换}：先做变换再变换一次，顺序不能换------所以
  \texttt{AB\ ≠\ BA} 一般情况下成立。
\end{enumerate}

\textbf{为什么这对科学计算重要？}
因为很多模型都是一句''线性组合''：多项式拟合是''基函数的线性组合''、线性回归是''特征的线性组合''、傅里叶分析是''正弦波的线性组合''。学会线性代数，等于看懂了这些模型共同的骨架。

\begin{quote}
\textbf{正文见}：\href{../../chapters/01-numpy/03-线性代数.md}{1 NumPy ·
03
线性代数}（矩阵运算与求解）、\href{../../chapters/01-numpy/04-多项式与拟合.md}{1
NumPy · 04 多项式与拟合}（最小二乘）。
\end{quote}

\begin{center}\rule{0.5\linewidth}{0.5pt}\end{center}

\section{A.2
关键概念讲解：把名词翻译成人话}\label{a.2-ux5173ux952eux6982ux5ff5ux8bb2ux89e3ux628aux540dux8bcdux7ffbux8bd1ux6210ux4ebaux8bdd}

下面每个概念，先给定义，再用一句话人话解释，最后配一个 2×2 小例子。

\subsection{A.2.1
转置、内积与范数}\label{a.2.1-ux8f6cux7f6eux5185ux79efux4e0eux8303ux6570}

{\def\LTcaptype{none} % do not increment counter
\begin{longtable}[]{@{}lll@{}}
\toprule\noalign{}
概念 & 定义 & 人话 \\
\midrule\noalign{}
\endhead
\bottomrule\noalign{}
\endlastfoot
转置 \texttt{A\^{}T} & 行列互换 & ``把表格行变列、列变行'' \\
内积 \texttt{x·y\ =\ x\^{}T\ y} & 对应元素相乘再求和 &
``两个向量有多'同向''' \\
范数 \texttt{\textbar{}\textbar{}x\textbar{}\textbar{}\_2} &
\texttt{sqrt(x\^{}T\ x)} & ``向量长度'' \\
\end{longtable}
}

例子：\texttt{x\ =\ (1,\ 2)}，\texttt{y\ =\ (3,\ 4)}，则
\texttt{x·y\ =\ 1×3\ +\ 2×4\ =\ 11}，\texttt{\textbar{}\textbar{}x\textbar{}\textbar{}\_2\ =\ sqrt(1+4)\ ≈\ 2.24}。内积越大且同向，相似度越高------这就是余弦相似度、皮尔逊相关、注意力机制的底层直觉。

\subsection{A.2.2
逆、行列式与秩：能否把变换''倒回去''}\label{a.2.2-ux9006ux884cux5217ux5f0fux4e0eux79e9ux80fdux5426ux628aux53d8ux6362ux5012ux56deux53bb}

\begin{itemize}
\tightlist
\item
  \textbf{逆
  \texttt{A\^{}\{-1\}}}：把变换倒着做一遍。\texttt{A\^{}\{-1\}A\ =\ I}。
\item
  \textbf{行列式
  \texttt{det(A)}}：变换后''体积''放大的倍数。\texttt{det(A)\ =\ 0}
  表示变换把高维压扁了（不可逆），方程组要么无解要么无穷多解。
\item
  \textbf{秩
  \texttt{rank(A)}}：变换后空间的真正维度（列向量中线性无关的个数）。满秩方阵才可逆。
\end{itemize}

例子：\texttt{A\ =\ {[}{[}2,1{]},{[}1,2{]}{]}}。\texttt{det(A)\ =\ 4-1\ =\ 3\ ≠\ 0}，可逆；\texttt{A\^{}\{-1\}\ =\ (1/3){[}{[}2,-1{]},{[}-1,2{]}{]}}。而
\texttt{B\ =\ {[}{[}1,2{]},{[}2,4{]}{]}}
两列成比例，\texttt{rank(B)\ =\ 1}，\texttt{det(B)\ =\ 0}------它把平面压成一条线，因此
\texttt{Bx\ =\ b} 要么无解要么无穷多解。

\subsection{A.2.3
特征值与特征向量：变换的''不动轴''}\label{a.2.3-ux7279ux5f81ux503cux4e0eux7279ux5f81ux5411ux91cfux53d8ux6362ux7684ux4e0dux52a8ux8f74}

定义：若 \texttt{Av\ =\ λv}（\texttt{v\ ≠\ 0}），则 \texttt{v}
是特征向量，\texttt{λ}
是对应特征值。人话：\textbf{有的方向被矩阵只''拉伸''不''拐弯''，拉伸倍数是特征值}。

继续用 \texttt{A\ =\ {[}{[}2,1{]},{[}1,2{]}{]}}：特征值
\texttt{λ1\ =\ 3}（方向 \texttt{(1,1)}）、\texttt{λ2\ =\ 1}（方向
\texttt{(1,-1)}）。也就是说：沿对角线方向放 3 倍、沿反对角线方向不动。

用处： - 对称矩阵
\texttt{A\ =\ QΛQ\^{}T}：把坐标轴转成''特征方向''，问题立刻解耦（PCA
就是找这些方向）； -
幂迭代、PageRank、马尔可夫链稳态都在找''最大特征值的特征向量''。

\subsection{A.2.4
条件数：解的''敏感度''}\label{a.2.4-ux6761ux4ef6ux6570ux89e3ux7684ux654fux611fux5ea6}

\texttt{cond(A)\ =\ \textbar{}\textbar{}A\textbar{}\textbar{}\ ·\ \textbar{}\textbar{}A\^{}\{-1\}\textbar{}\textbar{}}，粗略说明：\texttt{b}
变化一点点，\texttt{x} 可能变化多大倍数。条件数大（病态）≠
无解，而是\textbf{数值上很难算准}。

\begin{quote}
\textbf{正文见}：\href{../../chapters/01-numpy/03-线性代数.md}{1 NumPy ·
03
线性代数}（条件数/病态）、\href{../../chapters/03-scipy/02-优化工具包.md}{3
SciPy · 02 优化工具包}（优化里的病态）。
\end{quote}

\begin{center}\rule{0.5\linewidth}{0.5pt}\end{center}

\section{A.3
分解与算法：为什么这样做}\label{a.3-ux5206ux89e3ux4e0eux7b97ux6cd5ux4e3aux4ec0ux4e48ux8fd9ux6837ux505a}

\subsection{A.3.1 解方程组：从高斯消元到
LU}\label{a.3.1-ux89e3ux65b9ux7a0bux7ec4ux4eceux9ad8ux65afux6d88ux5143ux5230-lu}

求解 \texttt{Ax\ =\ b}
最直接是高斯消元：把增广矩阵化成上三角，再回代。消元过程等价于把
\texttt{A} 写成 \texttt{A\ =\ LU}（L 下三角、U
上三角）。好处是：\textbf{解多组 \texttt{b} 时复用
LU，一次消元多次求解}。

\begin{itemize}
\tightlist
\item
  矩阵规模 \texttt{n×n}，复杂度 \texttt{O(n\^{}3)}；
\item
  数值上要\textbf{部分主元}（选最大的元素当主元），避免小除数放大误差；
\item
  Python 里直接 \texttt{np.linalg.solve}（内部就是带主元的
  LU），\textbf{不要}用 \texttt{inv(A)\ @\ b}。
\end{itemize}

\subsection{A.3.2 特征值分解与
SVD：谁是''最重要的方向''}\label{a.3.2-ux7279ux5f81ux503cux5206ux89e3ux4e0e-svdux8c01ux662fux6700ux91cdux8981ux7684ux65b9ux5411}

\begin{itemize}
\tightlist
\item
  \textbf{特征分解}：\texttt{A\ =\ VΛV\^{}\{-1\}}（对称阵时 \texttt{V}
  正交，\texttt{A\ =\ QΛQ\^{}T}）。适合方阵、可对角化矩阵；求
  \texttt{A\^{}k}、判断稳态很舒服。
\item
  \textbf{奇异值分解
  SVD}：\texttt{A\ =\ UΣV\^{}T}，\textbf{任意矩阵都行}。\texttt{U}、\texttt{V}
  正交，\texttt{Σ}
  对角（奇异值按降序）。直觉：把''变换''分解成''先旋转、再缩放、再旋转''。
\end{itemize}

SVD 最漂亮的性质是\textbf{最优低秩近似}（Eckart--Young）：保留前
\texttt{k} 大奇异值，就得到与原矩阵''最接近''的秩 \texttt{k}
矩阵。所以： - 图像压缩：把大矩阵拆成''少数奇异值 × 两个小矩阵''； -
PCA：对数据矩阵做 SVD，右奇异向量就是主成分方向； -
伪逆与最小二乘：\texttt{A\^{}+\ =\ VΣ\^{}\{-1\}U\^{}T}，比正规方程稳定。

\begin{quote}
\textbf{正文见}：\href{../../chapters/01-numpy/05-综合案例.md}{1 NumPy ·
05 综合案例}（SVD
图像压缩）、\href{../../chapters/08-sklearn/03-无监督学习的案例.md}{8
sklearn · 03 无监督学习的案例}（PCA 降维）。
\end{quote}

\subsection{A.3.3
最小二乘：数据有噪声怎么办}\label{a.3.3-ux6700ux5c0fux4e8cux4e58ux6570ux636eux6709ux566aux58f0ux600eux4e48ux529e}

超定方程
\texttt{Ax\ ≈\ b}（方程比未知数多）通常无精确解，最小二乘找''残差平方和最小''的解：

\begin{Shaded}
\begin{Highlighting}[]
\NormalTok{正规方程：x\^{} = (A\^{}T A)\^{}\{{-}1\} A\^{}T b}
\NormalTok{几何：    把 b 投影到 A 的列空间，误差向量垂直于列空间}
\NormalTok{数值：    优先 lstsq（QR/ SVD 路径），避免求 (A\^{}T A)\^{}\{{-}1\}（会平方条件数）}
\end{Highlighting}
\end{Shaded}

如果再加 \texttt{\textbar{}\textbar{}x\textbar{}\textbar{}\^{}2}
惩罚（岭回归），就是''在最小二乘和'不放大'之间妥协''------这是第 8
章正则化的数学原型。

\begin{center}\rule{0.5\linewidth}{0.5pt}\end{center}

\section{A.4 Python
对应（速查）}\label{a.4-python-ux5bf9ux5e94ux901fux67e5}

\begin{Shaded}
\begin{Highlighting}[]
\ImportTok{import}\NormalTok{ numpy }\ImportTok{as}\NormalTok{ np, sympy }\ImportTok{as}\NormalTok{ sp}

\NormalTok{A }\OperatorTok{=}\NormalTok{ np.array([[}\DecValTok{2}\NormalTok{, }\DecValTok{1}\NormalTok{], [}\DecValTok{1}\NormalTok{, }\DecValTok{2}\NormalTok{]])}
\NormalTok{b }\OperatorTok{=}\NormalTok{ np.array([[}\DecValTok{1}\NormalTok{], [}\DecValTok{2}\NormalTok{]])}

\NormalTok{np.linalg.solve(A, b)                 }\CommentTok{\# 解方程组（推荐）}
\NormalTok{np.linalg.inv(A)                      }\CommentTok{\# 显式求逆（尽量不用）}
\NormalTok{np.linalg.det(A), np.linalg.matrix\_rank(A)}
\NormalTok{np.linalg.eig(A)                      }\CommentTok{\# 特征值/特征向量}
\NormalTok{np.linalg.eigh(A) }\ControlFlowTok{if}\NormalTok{ (A }\OperatorTok{==}\NormalTok{ A.T).}\BuiltInTok{all}\NormalTok{() }\ControlFlowTok{else} \VariableTok{None}   \CommentTok{\# 对称矩阵专用}
\NormalTok{np.linalg.svd(A, full\_matrices}\OperatorTok{=}\VariableTok{False}\NormalTok{)  }\CommentTok{\# U, S, Vt（S 是向量!）}
\NormalTok{np.linalg.lstsq(A, b, rcond}\OperatorTok{=}\VariableTok{None}\NormalTok{)      }\CommentTok{\# 最小二乘}
\NormalTok{np.linalg.cond(A)                      }\CommentTok{\# 条件数}
\NormalTok{np.linalg.pinv(A)                      }\CommentTok{\# 伪逆}

\CommentTok{\# SymPy 符号解}
\NormalTok{sp.Matrix([[}\DecValTok{2}\NormalTok{,}\DecValTok{1}\NormalTok{],[}\DecValTok{1}\NormalTok{,}\DecValTok{2}\NormalTok{]]).eigenvects()}
\end{Highlighting}
\end{Shaded}

{\def\LTcaptype{none} % do not increment counter
\begin{longtable}[]{@{}ll@{}}
\toprule\noalign{}
你想要 & 用哪个 \\
\midrule\noalign{}
\endhead
\bottomrule\noalign{}
\endlastfoot
解方阵方程组 & \texttt{np.linalg.solve} \\
最小二乘 & \texttt{np.linalg.lstsq} \\
特征值（对称） & \texttt{np.linalg.eigh} \\
SVD & \texttt{np.linalg.svd}（注意 \texttt{S} 是向量） \\
判断病态 & \texttt{np.linalg.cond} \\
\end{longtable}
}

\begin{center}\rule{0.5\linewidth}{0.5pt}\end{center}

\section{A.5 动手例题（选做：手算 +
验证）}\label{a.5-ux52a8ux624bux4f8bux9898ux9009ux505aux624bux7b97-ux9a8cux8bc1}

\textbf{例 1：解方程组}。解 @BT@2x + y = 1@BT@、@BT@x + 2y =
2@BT@。手算：由第一式 @BT@y = 1 - 2x@BT@ 代入第二式 @BT@x + 2(1-2x) =
2@BT@ → @BT@-3x = 0@BT@ → @BT@x = 0@BT@，@BT@y =
1@BT@。代码验证：@BT@np.linalg.solve({[}{[}2,1{]},{[}1,2{]}{]},
{[}1,2{]}) = {[}0, 1{]}@BT@。

\textbf{例 2：最小二乘的''投影直觉''}。数据 @BT@(0,1), (1,2), (2,3)@BT@
明显在同一直线 @BT@y = x + 1@BT@ 上，最小二乘应还原出斜率 1、截距
1。若把第 3 点改成 @BT@(2, 4)@BT@（有噪声），斜率会略大于
1------因为解被''拉''向那个异常点。试试：

\begin{Shaded}
\begin{Highlighting}[]
\ImportTok{import}\NormalTok{ numpy }\ImportTok{as}\NormalTok{ np}
\NormalTok{x }\OperatorTok{=}\NormalTok{ np.array([}\FloatTok{0.}\NormalTok{, }\FloatTok{1.}\NormalTok{, }\FloatTok{2.}\NormalTok{])}\OperatorTok{;}\NormalTok{ y }\OperatorTok{=}\NormalTok{ np.array([}\FloatTok{1.}\NormalTok{, }\FloatTok{2.}\NormalTok{, }\FloatTok{3.}\NormalTok{])}
\NormalTok{A }\OperatorTok{=}\NormalTok{ np.vstack([x, np.ones\_like(x)]).T        }\CommentTok{\# 列：[x, 1]}
\NormalTok{k, b }\OperatorTok{=}\NormalTok{ np.linalg.lstsq(A, y, rcond}\OperatorTok{=}\VariableTok{None}\NormalTok{)[}\DecValTok{0}\NormalTok{]}
\BuiltInTok{print}\NormalTok{(k, b)                                   }\CommentTok{\# 1.0 1.0}
\NormalTok{y2 }\OperatorTok{=}\NormalTok{ np.array([}\FloatTok{1.}\NormalTok{, }\FloatTok{2.}\NormalTok{, }\FloatTok{4.}\NormalTok{])}
\BuiltInTok{print}\NormalTok{(np.linalg.lstsq(A, y2, rcond}\OperatorTok{=}\VariableTok{None}\NormalTok{)[}\DecValTok{0}\NormalTok{])  }\CommentTok{\# 斜率 \textgreater{} 1，被噪声拉偏}
\end{Highlighting}
\end{Shaded}

\section{A.6 常见误区 +
用在哪（双向）}\label{a.6-ux5e38ux89c1ux8befux533a-ux7528ux5728ux54eaux53ccux5411}

{\def\LTcaptype{none} % do not increment counter
\begin{longtable}[]{@{}ll@{}}
\toprule\noalign{}
误区 & 正确 \\
\midrule\noalign{}
\endhead
\bottomrule\noalign{}
\endlastfoot
\texttt{A\ *\ B} 当矩阵乘 & \texttt{*} 是逐元素；用 \texttt{A\ @\ B} \\
用 \texttt{inv(A)\ @\ b} & 用 \texttt{solve} / \texttt{lstsq} \\
以为 \texttt{S} 是矩阵 & SVD 返回奇异值\textbf{向量}，需
\texttt{np.diag(S)} \\
忽略条件数 & 病态时解会''抖''，先看 \texttt{cond} \\
非对称矩阵用 \texttt{eigh} & \texttt{eigh} 要求对称；否则用
\texttt{eig} \\
\end{longtable}
}

\textbf{使用章节（可点击）} \textbar{} 章 \textbar{} 哪里用到 \textbar{}
链接 \textbar{} \textbar{} ---- \textbar{} ---- \textbar{} ----
\textbar{} \textbar{} 1 NumPy \textbar{} 03 线性代数、05 SVD 压缩
\textbar{} \href{../../chapters/01-numpy/03-线性代数.md}{正文}
\textbar{} \textbar{} 2 SymPy \textbar{} 矩阵与特征值 \textbar{}
\href{../../chapters/02-sympy/01-符号对象与基本运算.md}{正文} \textbar{}
\textbar{} 3 SciPy \textbar{} 优化/插值中的矩阵 \textbar{}
\href{../../chapters/03-scipy/02-优化工具包.md}{正文} \textbar{}
\textbar{} 8 sklearn \textbar{} PCA/降维 \textbar{}
\href{../../chapters/08-sklearn/03-无监督学习的案例.md}{正文} \textbar{}

\textbf{下游衔接}：intro-mathmodel 第 1 章（Numpy 与线性代数）、第 3
章（线性规划→非线性规划）。 \textbf{延伸阅读}：mml-book 第
2\textasciitilde4 章、SciPy Lectures、numpy.linalg 官方文档（见
\href{./references.md}{本目录 references.md}）。

\begin{center}\rule{0.5\linewidth}{0.5pt}\end{center}

\section{A.7 补充：特征值与 SVD 的 3
个实际用途}\label{a.7-ux8865ux5145ux7279ux5f81ux503cux4e0e-svd-ux7684-3-ux4e2aux5b9eux9645ux7528ux9014}

\begin{enumerate}
\def\labelenumi{\arabic{enumi}.}
\tightlist
\item
  \textbf{收敛性判断（马尔可夫链/PageRank）}：状态转移矩阵的最大特征值决定系统是否收敛到稳态；PageRank
  就是求转移矩阵最大特征值的特征向量（幂迭代）。
\item
  \textbf{去噪与压缩（SVD）}：保留前 k 个奇异值 =
  丢掉''不重要的方向''。对一张 100×100 的图，只保留前 10
  个奇异值，视觉上几乎一样，但数据量从 10000 降到约 2100。
\item
  \textbf{求逆与最小二乘（伪逆）}：@BT@A\^{}+ =
  VΣ\textsuperscript{\{-1\}U}T@BT@ 对任意矩阵（含奇异）都定义良好，是
  @BT@lstsq@BT@ 的底层路径。
\end{enumerate}

\textbf{一句话总结}：特征值/奇异值大小 =
该方向的重要程度；``取前几个''就是''保留主要信息，丢弃噪声''。
